\section{De RLHF a DPO}
\subsection{Las dos etapas de RLHF}
\begin{enumerate}
  \item Recopilar datos de preferencia $\to$ entrenar el Reward Model
  \item Usar el RM para construir una pérdida sustituta $\to$ optimizar la política con PPO
\end{enumerate}

\subsection{Derivación de DPO}
El conocimiento clave de DPO: existe una relación de forma cerrada entre la política óptima $\pi^*$ y la reward function. Aprovechando esta relación, se puede optimizar la política directamente a partir de los datos de preferencia, sin un RM explícito.

$$\mathcal{L}_{\text{DPO}} = -\mathbb{E}\left[\log \sigma\left(\beta \log \frac{\pi_\theta(y_w|x)}{\pi_{\text{ref}}(y_w|x)} - \beta \log \frac{\pi_\theta(y_l|x)}{\pi_{\text{ref}}(y_l|x)}\right)\right]$$

\subsection{Resumen de la sección}
DPO comprime, mediante una derivación matemática, las dos etapas de RLHF en un solo paso: una simplificación elegante.

\newpage
